\section{Control and Protocol}
\label{sec:proto}

The host drives the fabric through a byte-oriented UART protocol interpreted by
a control finite-state machine (FSM). The protocol is deliberately minimal: six
single-byte commands, each with a fixed payload and reply, so that the same
message stream can be produced and checked by both the hardware and a software
model.

\subsection{Command set}

Table~\ref{tab:proto} lists the commands. Configuration and input transfers
carry an 8-bit additive checksum; the device rejects a corrupted transfer with
a negative acknowledgement, and appends a checksum to the register read-back.
The identification reply reports the array geometry, so the host can adapt at
run time rather than assuming fixed dimensions.

\begin{table}[t]
  \centering
  \caption{UART command set (protocol version 2).}
  \label{tab:proto}
  \small
  \begin{tabular}{@{}llp{4.6cm}@{}}
    \toprule
    Cmd & Payload & Reply \\
    \midrule
    \cmd{ID}  & --- & signature, version, ROWS, COLS, width \\
    \cmd{CFG} & 16 words $+$ cksum & ACK / NACK \\
    \cmd{WR}  & 8 words $+$ cksum  & ACK / NACK \\
    \cmd{RUN} & step count $N$     & ACK (after $N$ steps) \\
    \cmd{RD}  & ---                & 16 registers $+$ cksum \\
    \cmd{RST} & ---                & ACK (clears registers) \\
    \bottomrule
  \end{tabular}
\end{table}

\subsection{A transaction on the wire}

To make the framing concrete, consider writing the edge inputs
$\text{west}=(1,0,0,0)$ and $\text{north}=(5,0,0,0)$ with the \cmd{WR} command.
The host sends the command byte \byte{03} followed by the four west then four
north values as 16-bit little-endian words, and finally the additive checksum:

\begin{center}\small\ttfamily
03\quad 0100 0000 0000 0000\quad 0500 0000 0000 0000\quad 06
\end{center}

\noindent The sixteen payload bytes sum to $1+5=6$, so the trailing checksum is
\byte{06}; the device recomputes the same sum, finds it matches, and replies
with a single \cmd{ACK} byte (\byte{79}). A corrupted payload would sum
differently and draw a \cmd{NACK} (\byte{1F}) instead, on which the host
retransmits. Every multi-byte transfer follows this shape---command, fixed
payload, one checksum byte, one-byte acknowledgement---which is what lets the
same stream be produced and checked identically by the hardware FSM and the
software model.

\subsection{Controller state machine}

Figure~\ref{fig:fsm} shows the control FSM. From the idle state the first
received byte selects a branch. The short transactions (\cmd{ID}, \cmd{RST} and
the rejection of an unknown command) reply immediately through a shared
transmit handshake. The bulk transfers accumulate their payload, verify (or
emit) a checksum in a dedicated check state, and then reply. The \cmd{RUN}
branch asserts \emph{step} once per remaining cycle before acknowledging.

\begin{figure*}[t]
  \centering
  % figures/fsm.tex — controller finite-state machine (protocol v2, logical).
\begin{tikzpicture}[->, >={Latex[length=2mm]}, x=17mm, y=13mm,
                    every state/.style={st}, font=\scriptsize,
                    every edge/.style={draw, ->, >={Latex[length=2mm]}}]
  % states (absolute placement for a clean left-to-right flow)
  \node[st, fill=cStateE, text=white] (idle) at (0,0)    {IDLE};
  \node[st] (cfg)   at (2, 2)    {CFG};
  \node[st] (cfgck) at (4, 2)    {CFG\\CK};
  \node[st] (wr)    at (2, 0.8)  {WR};
  \node[st] (wrck)  at (4, 0.8)  {WR\\CK};
  \node[st] (run)   at (2,-0.6)  {RUN};
  \node[st] (rd)    at (2,-2)     {RD};
  \node[st] (rdck)  at (4,-2)     {RD\\CK};
  \node[st] (send)  at (6, 0)     {SEND};

  % decode: idle -> command branches
  \path (idle) edge node[above,sloped]{CFG} (cfg)
        (idle) edge node[above,sloped]{WR}  (wr)
        (idle) edge node[below,sloped]{RUN} (run)
        (idle) edge node[below,sloped]{RD}  (rd);

  % payload accumulate + checksum
  \path (cfg) edge node[above]{16 words} (cfgck)
        (wr)  edge node[above]{8 words}  (wrck)
        (rd)  edge node[above]{16 regs}  (rdck);

  % replies converge on the shared SEND collector
  \path (cfgck) edge node[above,sloped]{ACK/NACK} (send)
        (wrck)  edge node[above,sloped]{ACK/NACK} (send)
        (rdck)  edge node[below,sloped]{$+$cksum}  (send)
        (run)   edge node[above,sloped]{$N{=}0$: ACK} (send);
  \path (run) edge[loop below] node{step, $N{-}1$} (run);

  % short commands reply directly; SEND returns to IDLE
  \path (idle) edge[bend right=55] node[below,sloped]{ID / RST / bad} (send);
  \path (send) edge[bend right=40] node[above]{tx done} (idle);

  % watchdog abort on stalled payload
  \path[dashed] (cfg) edge[bend right=18] node[above,sloped]{\itshape timeout} (idle);
\end{tikzpicture}

  \caption{Control finite-state machine. Payload-bearing commands (\cmd{CFG},
  \cmd{WR}, \cmd{RD}) pass through a checksum state before the shared send
  handshake; \cmd{RUN} pulses \emph{step} $N$ times. A watchdog (dashed)
  returns the machine to idle if a payload stalls, so a desynchronised host
  always recovers.}
  \label{fig:fsm}
\end{figure*}

A watchdog timer guards every state that waits for payload bytes: if no byte
arrives within a bounded interval the machine aborts back to idle. Together with
the checksum, this makes the link self-recovering---a property the host side
mirrors by retransmitting rejected transfers and resynchronising after a failed
run.
